%% ===================== NEW FRONT MATTER (consolidated paper) =====================
\title[Correlations of Artin status among primes]{Correlations of Artin status among primes:\\ consecutive-prime anticorrelation, cross-base entanglement,\\ and an elementary decomposition}

\author{Josh Bald}
\address{Ontario, Canada}
\email{jpbald93@gmail.com}
\thanks{\textbf{Corresponding author.} Josh Bald, Independent Researcher.
	Email: \texttt{jpbald93@gmail.com}. ORCID: \texttt{0009-0002-1317-6489}.}

\subjclass[2020]{Primary 11A07; Secondary 11N05, 11N13, 11Y16, 11Y60}
\keywords{Artin's primitive root conjecture, primitive roots, consecutive primes,
	Lemke Oliver--Soundararajan bias, entanglement, quadratic characters,
	Legendre symbol, full reptend primes, computational number theory}

\date{September 2026}

\begin{abstract}
	Say a prime $p$ is \emph{Artin base $a$} if $a$ is a primitive root modulo $p$.
	This paper studies correlations of Artin status along two axes --- \emph{consecutive
	primes} at a fixed base, and \emph{several bases} at a fixed prime --- for all primes up to $10^9$ (with $50{,}847{,}531$ of them, $p\ge 7$, in the
	consecutive census), and then asks what, if anything, such correlations are
	computationally worth.

	Along the first axis, Artin statuses of consecutive primes are measurably
	\emph{anticorrelated} at base $10$: $P(\Art_{n+1}\mid\Art_n)=0.36514$ against
	$P(\Art_{n+1}\mid\lnot\Art_n)=0.37928$, a deficit $\delta=-0.01414$
	(Pearson $10{,}167$ on one degree of freedom; signed square root $-100.8$ under a
	nominal independent-status reference, over $50{,}847{,}530$ pairs). The dependence
	has strong gap structure and contains one deterministic component, an elementary
	\emph{exclusion law}: if $p,p+g>5$ are prime with $g\equiv 20 \pmod{40}$, then $10$
	is a quadratic residue modulo exactly one of them, so the two can never both be
	Artin; among $2{,}195{,}882$ pairs at gaps $20$ and $60$ there is no doubly-Artin
	pair, and an independent recomputation extends this to $194{,}296{,}748$ such pairs
	below $10^{11}$, again with none. Conversely $g \equiv 0 \pmod{40}$ forces equal
	quadratic-residue status and strong positive dependence ($P(\Art_{n+1}\mid\Art_n)=0.899$
	at $g=40$). Conditioning on the joint residues of $(p_n,p_{n+1})$ modulo $120$ or
	$840$ places $98.7\%$ and $99.5\%$ of the global deficit between joint-residue cells
	rather than within them (an exact covariance decomposition; $96.1$--$99.2\%$ and $97.8$--$99.7\%$
	across the four weighting conventions we report). As an intermediate link we report the apparently unreported
	anticorrelation $r(\omega(p_n-1),\omega(p_{n+1}-1))=-0.041$.

	Along the second axis we record a second exclusion law: writing $d=\sqf(ab)$ for the
	squarefree part of $ab$, no prime with $\bigl(\tfrac{d}{p}\bigr)=-1$ is simultaneously
	Artin for $a$ and for $b$; whenever $a$ and $b$ have different squarefree parts this
	bars half of all primes, and it implies a triple form --- if $\sqf(c)=\sqf(ab)$ then no
	odd prime is Artin for $a,b,c$ at once (e.g.\ $2,5,10$). Across the $66$ base pairs of
	$\{2,3,5,6,7,10,11,13,15,17,21,29\}$ the measured same-prime correlation $\phi(a,b)$ is
	positive for $64$ pairs (mean $0.0352$); conditioning on $\omega(p-1)$ removes only
	$55\%$ of it, while conditioning on a fine signature of $p-1$ removes $95\%$. The
	residual is concentrated on the fifteen pairs whose product class completes a triple
	and lies in the quadratic-character channel. Joint densities predicted by the
	Matthews--Moree--Stevenhagen Kummer-degree heuristic match the measured densities to a
	mean relative error of $0.012\%$ with no fitted parameters.

	Finally we \emph{audit the computational content} of these correlations. In a resubstitution diagnostic the consecutive-prime correlation carries $0.000009$ nats
	($1.3\times10^{-5}$ bits) of conditional-entropy reduction beyond the elementary
	joint-residue baseline, which itself carries $0.4195$ nats. A simulation in which every
	status depends only on $p \bmod 840$ reproduces the part of that increment which exceeds
	the mod-$120$ baseline: residues (the prime $7$ of $p-1$) suffice to reproduce that excess
	without predecessor-specific information, and in held-out scoring the predecessor's status
	gives no gain beyond residues mod $840$.
Both exclusion laws are computationally redundant: the gap law from quadratic reciprocity
	together with multiplicativity, the same-prime law from multiplicativity alone. The quadratic filter expressed by these exclusion laws --- a primitive root must be a
quadratic non-residue --- is elementary, is already applied by a character-aware
implementation, is exactly what the observed conductor-$40$ channel expresses, and halves
the factorizations of $p-1$ in a scanning search; ordering
	the small prime channels of $p-1$ before the full factorization is about $13\%$ faster
	again.
We conclude that, for the consecutive-prime workloads tested, these correlations describe arithmetic structure
rather than supply algorithms, and we state what would change that verdict.
\end{abstract}

\maketitle

\section{Introduction}
\label{sec:intro}

Artin's conjecture on primitive roots, communicated to Hasse in 1927 (see
\cite{Artin1965}), concerns one base at a time: for non-square
$a \ne 0, \pm 1$ the set of primes $p$ for which $a$ is a primitive root mod $p$
should have positive density, a multiple of Artin's constant
$C = \prod_q \bigl(1 - \tfrac{1}{q(q-1)}\bigr) \approx 0.3739558$; for base $10$ no
correction factor applies, so the conjectural density is exactly $C$. Hooley~\cite{Hooley1967}
proved this under GRH; unconditionally, Gupta--Murty~\cite{GuptaMurty1984} and
Heath-Brown~\cite{HeathBrown1986} showed the conjecture can fail for at most two prime
bases, and Klurman--Shparlinski--Ter\"av\"ainen~\cite{KlurmanShparlinskiTeravainen2025}
established the predicted asymptotic on average over almost all small bases; recent work
has sharpened the analytic understanding for individual primes, with Fan and Pollack
\cite{FanPollack2025} proving a version uniform in the base and Perucca and Shparlinski
\cite{PeruccaShparlinski2025} giving uniform upper and lower bounds for the Artin density.
Moree's survey~\cite{Moree2012} is the standard account. For $a=10$ the property has a classical
meaning: $10$ is a primitive root mod $p$ exactly when the decimal expansion of $1/p$ has
maximal period $p-1$ --- the \emph{full reptend} primes~\cite[pp.~166--171]{ConwayGuy1996}.

Three further strands meet here. The first is the bias discovered numerically by
Knapowski--Tur\'an, Ko and Ash--Beltis--Gross--Sinnott and quantitatively treated by
Lemke Oliver and Soundararajan~\cite{LOS2016}: consecutive primes in the same reduced
residue class modulo $q$ occur markedly less often than equidistribution predicts, a
phenomenon explained conjecturally by the Hardy--Littlewood $k$-tuple
conjectures~\cite{HardyLittlewood1923}. Since the Artin property of $p$ is governed by the
factorisation of $p-1$, and since residues of consecutive primes modulo small $q$ are
correlated by~\cite{LOS2016}, one expects \emph{some} induced dependence between the Artin
indicators of consecutive primes. The second is the joint behaviour \emph{across bases} at
a fixed prime, computed under GRH by Matthews~\cite{Matthews1976} and systematically by
Moree--Stevenhagen~\cite{MoreeStevenhagen2014}, with the qualitative message that
$\Acal_a$ and $\Acal_b$ are \emph{not} independent and that the deviation is governed by
entanglement of the associated Kummer--Galois extensions. The framework remains active:
Lenstra's Galois-theoretic setting~\cite{Lenstra1977}, the unified Artin-type treatment of
J\"arviniemi--Perucca--Sgobba~\cite{JarviniemiPeruccaSgobba2025}, unconditional results in
cases where GRH can be circumvented~\cite{Sgobba2025}, and number-field variants
\cite{Kimmel2024}. The third is the question of
which primes at a fixed shift can \emph{both} be Artin: Tinkov\'a--Waxman--Zindulka
\cite{TinkovaWaxmanZindulka2023} identify shift--base pairs admitting no such prime, and
our Theorem~\ref{thm:exclusion} is the base-$10$ instance of that phenomenon; we claim no
priority for it.

The present paper asks the \emph{statistical} question along both axes --- how Artin
statuses of consecutive primes are correlated on average, and how Artin statuses of several
bases are correlated at the same prime --- and then asks a question we have not seen posed
in this literature: \emph{what is such a correlation computationally worth?}\footnote{An
earlier version of parts of this work appeared as the preprints \emph{Correlations between
primitive root statuses of consecutive primes} (Zenodo,
\href{https://doi.org/10.5281/zenodo.22863946}{doi:10.5281/zenodo.22863946}) and
\emph{Cross-base correlations of Artin primes} (Zenodo,
\href{https://doi.org/10.5281/zenodo.22878204}{doi:10.5281/zenodo.22878204}). This paper
consolidates those two preprints and adds the audit of \S\ref{sec:content}; it supersedes
them, and the preprints are retained for the record.} The answer we give is a measured
negative, and we regard it as part of the result rather than a failure of it.

Two comparisons should be made explicit. Garcia--Kahoro--Luca~\cite{GarciaKahoroLuca2019}
(see also~\cite{GarciaLucaShiUdell2020}) compare the \emph{number} of primitive roots,
$\varphi(p-1)$ versus $\varphi(p+1)$, for twin pairs; Pollack~\cite{Pollack2014} and
Baker--Pollack~\cite{BakerPollack2016} prove existence of runs of consecutive primes each
having a prescribed primitive root. Those are existence results, produced by sieves that
select rare favourable configurations; they say nothing about the typical joint behaviour
along the sequence of primes. The questions studied here are complementary: the average
correlation, its decomposition, and its computational content.

\section{Summary of results}
\label{sec:summary}

Throughout, \emph{Artin prime} means Artin prime for base $10$ unless stated otherwise,
$p_n$ denotes the $n$-th prime, $\omega(m)$ the number of distinct prime factors of $m$,
and $\sqf(m)$ the squarefree part of $m$.

\begin{enumerate}
	\item \textbf{Consecutive-prime exclusion law} (Theorem~\ref{thm:exclusion}). If $p$ and
	$p+g$ are primes with $p,p+g>5$ and $g\equiv 20 \pmod{40}$, then
	$\Lsym{10}{p+g}=-\Lsym{10}{p}$; consequently at most one of the two is Artin base $10$.
	If instead $g\equiv 0 \pmod{40}$ then $\Lsym{10}{p+g}=\Lsym{10}{p}$.
	\item \textbf{Same-prime exclusion law} (Theorem~\ref{thm:pair}). With $d=\sqf(ab)$, no
	prime with $\Lsym{d}{p}=-1$ is Artin for both $a$ and $b$. If $\sqf(c)=\sqf(ab)$ then no
	odd prime is Artin for $a$, $b$ and $c$ simultaneously (Theorem~\ref{thm:triple}).
	\item \textbf{Consecutive-prime anticorrelation} (\S\ref{sec:global}). Across all
	$50{,}847{,}530$ consecutive pairs up to $10^9$,
	$\delta = P(\Art_{n+1}\mid\Art_n)-P(\Art_{n+1}\mid\lnot\Art_n) = -0.01414$, with
	Pearson $\chi^2 = 10{,}167$ on one degree of freedom.
	\item \textbf{Gap structure} (\S\ref{sec:gaps}). $\delta(g)$ ranges from $-0.592$ at
	$g=60$, where the theorem forces exclusion, to $+0.643$ at $g=40$.
	\item \textbf{Residue conditioning} (\S\ref{sec:decomposition}). Conditioning on the
	joint residues of $(p_n,p_{n+1})$ modulo $120$ and $840$ gives an exact covariance
	decomposition of the global deficit, placing $98.69\%$ and $99.55\%$ of it between cells
	(within-cell terms $-0.000186$ and $-0.000064$); across four weighting conventions the
	between share lies in $[96.1\%,99.2\%]$ and $[97.8\%,99.7\%]$.
	\item \textbf{$\omega$ repulsion} (\S\ref{sec:omega}).
	$r(\omega(p_n-1),\omega(p_{n+1}-1)) = -0.0410$.
	\item \textbf{Cross-base correlation} (\S\ref{sec:measurement}). Over the $66$ pairs of
	the base set, $\phi(a,b)$ is positive for $64$ pairs (mean $0.0352$); the only negative
	pairs are $(2,10)$ and $(3,15)$.
	\item \textbf{Decomposition of the cross-base correlation} (\S\ref{sec:decomposition_crossbase}).
	Conditioning on $\omega(p-1)$ removes $55\%$ of the mean signed correlation; conditioning
	on the fine signature $\mathrm{sig}(p)$ removes $95\%$ of it.
	\item \textbf{Residual localisation} (\S\ref{sec:entanglement}). The surviving residual
	is concentrated on the triple-completing pairs and lies in the quadratic-character
	channel.
	\item \textbf{Kummer-degree densities} (\S\ref{sec:model}). The
	Matthews--Moree--Stevenhagen predictions match the measured joint densities to a mean
	relative error of $0.012\%$, with no fitted parameters.
	\item \textbf{Computational content: a measured negative} (\S\ref{sec:content}). The
	consecutive-prime correlation adds $0.000009$ nats of in-sample conditional-entropy
	reduction to the mod-$120$ joint-residue baseline; a residue-only simulation mod $840$
	reproduces that excess without predecessor-specific information, and no held-out estimator we tested gains
	beyond residues mod $840$; both exclusion laws are computationally
	redundant, the same-prime law from character multiplicativity alone and the gap law from
	reciprocity together with it;
	the quadratic filter expressed by these exclusion laws is the elementary non-residue test,
	which a character-aware implementation already applies and the conductor-$40$ channel expresses.
\end{enumerate}

\section{Preliminaries and notation}
\label{sec:prelim}

For an odd prime $p$ not dividing $a$, write $\Lsym{a}{p}$ for the Legendre symbol and
$\chi_a(p)=\Lsym{a}{p}$ for the quadratic character of conductor $\mathrm{cond}(\chi_a)$.
Write $\Acal_a$ for the set of primes for which $a$ is a primitive root, and
$X_a(p)=\mathbf{1}[\,p\in\Acal_a\,]$ for its indicator; along the consecutive-prime axis we
abbreviate $\Art_n = \mathbf{1}[\,p_n\in\Acal_{10}\,]$ for Artin status at the $n$-th prime.
Then $a$ is a primitive root modulo $p$ exactly when $a^{(p-1)/\ell}\not\equiv 1 \pmod p$ for every prime
$\ell \mid p-1$. Two elementary facts are used repeatedly. First, if $X_a(p)=1$ then
$\Lsym{a}{p}=-1$: a primitive root has order $p-1$, so $a^{(p-1)/2}\equiv -1$, whereas a
quadratic residue satisfies $a^{(p-1)/2}\equiv 1$. Second, quadratic characters are
multiplicative in the numerator, $\Lsym{ab}{p}=\Lsym{a}{p}\Lsym{b}{p}$, and the Legendre
symbol depends on $p$ only through $p \bmod 4d$ when $d=\sqf(ab)$.

The Artin property of $p$ is determined by the prime factorisation of $p-1$: the criterion
consults $\ell=2$ and every prime divisor $\ell$ of $p-1$. This is why both decompositions
below are stratified by the \emph{fine signature}
\[
	\mathrm{sig}(p) = \Bigl(\min(v_2(p-1),4),\;\; \{\,q\le 13 : q\mid p-1\,\},\;\;
	\min(\omega_{>13}(p-1),3)\Bigr),
\]
a partition of the primes into at most $4\cdot 32\cdot 4 = 512$ cells (for odd $p$ one has
$v_2(p-1)\ge 1$, so the truncated valuation takes four values, not five). The signature
is a coarse summary of the part of $p-1$ that the criterion consults first for small bases;
it records neither all prime divisors of $p-1$ nor the power-residue outcomes the criterion
tests, and so does not determine Artin status.
